% Section 12.3, from lectures/L12/appendix/b_exercises.md. Each exercise names the part of the
% review it repeats, "(A.2)" in the lecture, as a first line: "Repeterar avsnitt 12.1.2".

\section{Övningar}
\label{c12:sec:exercises}\label{c12:app:b}

\begin{aside}
\textbf{Så kontrollerar du ditt arbete.} Övningarna är repetition av hela kursen, i samma stil som
provets frågor. Gör dem på papper, utan miniräknare, med referensbladet i bilaga~\ref{app:avr}
bredvid, och kontrollera mot lösningsförslagen i bilaga~\ref{app:answers}.
Övning~\ref{c12:ex:10}, kontrollen, görs först för hand och sedan i Microchip Studio.

Varje övning säger vilken del av \secref{c12:app:a} den repeterar. Går en övning dåligt, läs den
delen och det kapitel den pekar på innan du går vidare.

Varje övning är märkt med sin sort. Exakt en övning är en \textbf{kontrolluppgift}, och i det här
kapitlet är det övning~\ref{c12:ex:10}.
\end{aside}

\exercise{Talomvandlingar}{Räkna för hand}
\label{c12:ex:1}
\emph{Repeterar \secref{c12:a2}.}

Fyll i tabellen. Visa hur du räknar i minst en omvandling åt varje håll.

\begin{filltable}{@{}YYY@{}}
  \thead{Binärt} & \thead{Decimalt} & \thead{Hexadecimalt} \\
  \midrule
  \code{1011 0101} & & \fillrule
  & 200 & \fillrule
  & & \code{0x7E} \fillrule
  & 64 & \\
\end{filltable}

\exercise{Binär addition och flaggor}{Räkna för hand}
\label{c12:ex:2}
\emph{Repeterar \secref{c12:a2} och \secref{c12:a6}.}
\begin{parts}
  \item Addera \code{0110 1101} och \code{0101 1011} binärt, med minnessiffror. Kontrollera svaret
    decimalt.
  \item Additionen görs med \code{add} i ett 8-bitars register. Blir flaggan \code{C} satt? Blir
    \code{N} satt?
  \item Läs resultatet som ett tal \textbf{med tecken}, i tvåkomplement. Vilket tal är det? Varför
    är det ett konstigt svar på en addition av två positiva tal?
\end{parts}

\exercise{Koder}{Räkna för hand}
\label{c12:ex:3}
\emph{Repeterar \secref{c12:a2}.}
\begin{parts}
  \item Skriv 47 i BCD.
  \item Skriv de åtta 3-bitars Graykoderna i ordning, från \code{000}. Hur många bitar ändras mellan
    två koder som står bredvid varandra?
  \item Vilka ASCII-koder har tecknen i texten \code{OK}? Bokstäverna ligger i alfabetisk ordning
    från \code{'A'} = \code{0x41}.
\end{parts}

\exercise{Vilken grind är det?}{Räkna för hand}
\label{c12:ex:4}
\emph{Repeterar \secref{c12:a3}.}

\bookfigure{l12_review_network}{Ett grindnät med en NOR-, en AND- och en
OR-grind.}{c12:fig:review_network}

\begin{parts}
  \item Skriv uttrycket vid varje grinds utgång, och uttrycket för X.
  \item Gör sanningstabellen för X.
  \item Vilken enda grind gör samma sak som hela nätet?
\end{parts}

\exercise{Förenkla}{Räkna för hand}
\label{c12:ex:5}
\emph{Repeterar \secref{c12:a3}.}

Förenkla algebraiskt, och ange vilken lag varje steg använder. Kontrollera svaret med en
sanningstabell.
\begin{parts}
  \item \code{X = ABC + ABC' + AB'C}
  \item \code{X = (A + B)(A + B')}
\end{parts}

\exercise{Två NAND-grindar}{Konstruktion}
\label{c12:ex:6}
\emph{Repeterar \secref{c12:a3} och \secref{c12:a4}.}
\begin{parts}
  \item Visa med De Morgan att \code{X = AB + C'} kan skrivas \code{X = ((AB)' · C)'}.
  \item Rita nätet med enbart NAND-grindar med två ingångar. Hur många behövs?
  \item Hur många kapslar 74HC00 behövs, och hur många grindar blir över?
\end{parts}

\exercise{Ett kontaktnät}{Konstruktion}
\label{c12:ex:7}
\emph{Repeterar \secref{c12:a3}.}

Lampan H1 ska styras av tre tryckknappar enligt \code{H1 = S1(S2 + S3')}.
\begin{parts}
  \item Rita kontaktnätet mellan +24~V och 0~V. Vilka kontakter är slutande, och vilka är
    brytande?
  \item Gör sanningstabellen. För vilka kombinationer av nedtryckta knappar lyser lampan?
\end{parts}

\exercise{Karnaughdiagram}{Räkna för hand}
\label{c12:ex:8}
\emph{Repeterar \secref{c12:a4}.}

En funktion av A, B, C och D är 1 för mintermerna 0, 1, 4, 5, 10, 11, 14 och 15, och 0 annars.
\begin{parts}
  \item Rita Karnaughdiagrammet, med \code{AB} på raderna och \code{CD} på kolumnerna, och fyll i
    ettorna.
  \item Gruppera minimalt, och skriv det minimerade uttrycket.
  \item Uttrycket är en känd grind med två av variablerna som ingångar. Vilken?
\end{parts}

\exercise{Mikrodatorn}{Förståelse}
\label{c12:ex:9}
\emph{Repeterar \secref{c12:a5}.}
\begin{parts}
  \item Rita mikrodatorns blockschema, med CPU (ALU, register, styrenhet, programräknare),
    programminne, dataminne, in- och utportar och bussar.
  \item Beskriv vad programräknaren gör under hämta-avkoda-utföra för en instruktion.
  \item ATmega328P startas om efter ett strömavbrott. Vad finns kvar i flashminnet, i SRAM och i
    registren?
\end{parts}

\exercise{Kontroll: följ programmet}{Kontroll}
\label{c12:ex:10}
\emph{Repeterar \secref{c12:a6}. Förutsäg för hand, kontrollera i simulatorn, förklara
skillnaden.}

\begin{asmcode}
.include "m328Pdef.inc"

.org 0x0000
    rjmp main

main:
    ldi r16, 0x9C
    ldi r17, 0x64
    add r16, r17
    ldi r18, 5
    sub r18, r17
    mov r19, r18
    lsr r19
    andi r19, 0x0F
end:
    rjmp end
\end{asmcode}

\task{a) För hand}
Fyll i tabellen med värdena \textbf{efter} varje instruktion, hexadecimalt. Fyll i flaggorna på
varje rad, men lämna en registerruta tom om registret inte har ändrats sedan raden ovanför.

\begin{filltable}{@{}lYYYYYYY@{}}
  \thead{Instruktion} & \thead{\code{r16}} & \thead{\code{r17}} & \thead{\code{r18}} &
    \thead{\code{r19}} & \thead{\code{C}} & \thead{\code{Z}} & \thead{\code{N}} \\
  \midrule
  \code{add r16, r17} & & & & & & & \fillrule
  \code{sub r18, r17} & & & & & & & \fillrule
  \code{lsr r19} & & & & & & & \fillrule
  \code{andi r19, 0x0F} & & & & & & & \\
\end{filltable}

\task{b) I simulatorn}
Skriv in programmet, stega det, och jämför med tabellen efter varje steg.

\task{c) Jämför}
Vilka rutor hade du fel i? Den vanligaste avvikelsen gäller \code{C} efter \code{andi}: vad visar
simulatorn, och varför?

\exercise{En loop}{Räkna för hand}
\label{c12:ex:11}
\emph{Repeterar \secref{c12:a6}.}

\begin{asmcode}
    ldi r20, 50
loop:
    nop
    dec r20
    brne loop
\end{asmcode}

\begin{parts}
  \item Hur många cykler tar koden, från \code{ldi} till och med den sista \code{brne}?
  \item Hur lång tid är det vid 16~MHz?
  \item Vilket värde ska laddas i \code{r20} för att koden ska ta exakt 25~µs?
\end{parts}

\exercise{Om, annars}{Program}
\label{c12:ex:12}
\emph{Repeterar \secref{c12:a6}.}

Lysdioden på PB0 ska lysa om värdet i \code{r16} är \textbf{minst 100}, och lysdioden på PB1
annars.
\begin{parts}
  \item Rita flödesplanen.
  \item Skriv programmet. Vilket villkorligt hopp används, och varför just det?
  \item Testa i simulatorn med \code{r16} = 150, 100 och 99. Vilket av de tre värdena är viktigast
    att testa, och varför?
\end{parts}

\exercise{Stacken}{Räkna för hand, Fördjupning}
\label{c12:ex:13}
\emph{Repeterar \secref{c12:a6}.}

\code{SP} är \code{0x08FF}. Programmet kör \code{rcall work}, och subrutinen \code{work} börjar med
\code{push r16} och \code{push r17}.
\begin{parts}
  \item Vad är \code{SP} efter de två \code{push}?
  \item Vilka byte ligger på adresserna \code{0x08FF}--\code{0x08FC}, i ord?
  \item Subrutinen slutar med \code{pop r16}, \code{pop r17} och \code{ret}. Programmet återvänder
    rätt, men något annat är fel. Vad?
\end{parts}
